\documentclass[10pt,a4paper]{amsart}
\usepackage{amsmath,amssymb,amsthm,graphicx}
% Retained minimal preamble for the Stacks Project chapter more-algebra.tex (More on Algebra).
% The chapter file pinned-original-source.tex carries only \input{preamble} on its first line, so no
% preamble text lives inside the chapter itself.  The Stacks Project uses one single shared preamble for
% every chapter; the snapshot stages it identically as sources/stacks-chow/preamble.tex and
% sources/stacks-etale-cohomology/preamble.tex.  The block below is copied mechanically, line for line,
% from sources/stacks-chow/preamble.tex lines 14-18 (the Xy-pic setup) and 154-217 (the native theorem
% declarations together with the mathematical macros).  Nothing was added, removed or redefined; the
% wrapper class and the amsmath/amssymb/amsthm/graphicx base packages are supplied by the assembler.
% Chapter: more-algebra.tex -- More on Algebra.
% For commutative diagrams we use Xy-pic
\usepackage[all]{xy}

% We use 2cell for 2-commutative diagrams.
\xyoption{2cell}

% Theorem environments.
%
\theoremstyle{plain}
\newtheorem{theorem}[subsection]{Theorem}
\newtheorem{proposition}[subsection]{Proposition}
\newtheorem{lemma}[subsection]{Lemma}

\theoremstyle{definition}
\newtheorem{definition}[subsection]{Definition}
\newtheorem{example}[subsection]{Example}
\newtheorem{exercise}[subsection]{Exercise}
\newtheorem{situation}[subsection]{Situation}

\theoremstyle{remark}
\newtheorem{remark}[subsection]{Remark}
\newtheorem{remarks}[subsection]{Remarks}

\numberwithin{equation}{subsection}

% Macros
%
\def\lim{\mathop{\mathrm{lim}}\nolimits}
\def\colim{\mathop{\mathrm{colim}}\nolimits}
\def\Spec{\mathop{\mathrm{Spec}}}
\def\Hom{\mathop{\mathrm{Hom}}\nolimits}
\def\Ext{\mathop{\mathrm{Ext}}\nolimits}
\def\SheafHom{\mathop{\mathcal{H}\!\mathit{om}}\nolimits}
\def\SheafExt{\mathop{\mathcal{E}\!\mathit{xt}}\nolimits}
\def\Sch{\mathit{Sch}}
\def\Mor{\mathop{\mathrm{Mor}}\nolimits}
\def\Ob{\mathop{\mathrm{Ob}}\nolimits}
\def\Sh{\mathop{\mathit{Sh}}\nolimits}
\def\NL{\mathop{N\!L}\nolimits}
\def\CH{\mathop{\mathrm{CH}}\nolimits}
\def\proetale{{pro\text{-}\acute{e}tale}}
\def\etale{{\acute{e}tale}}
\def\QCoh{\mathit{QCoh}}
\def\Ker{\mathop{\mathrm{Ker}}}
\def\Im{\mathop{\mathrm{Im}}}
\def\Coker{\mathop{\mathrm{Coker}}}
\def\Coim{\mathop{\mathrm{Coim}}}

% Boxtimes
%
\DeclareMathSymbol{\boxtimes}{\mathbin}{AMSa}{"02}

%
% Macros for moduli stacks/spaces
%
\def\QCohstack{\mathcal{QC}\!\mathit{oh}}
\def\Cohstack{\mathcal{C}\!\mathit{oh}}
\def\Spacesstack{\mathcal{S}\!\mathit{paces}}
\def\Quotfunctor{\mathrm{Quot}}
\def\Hilbfunctor{\mathrm{Hilb}}
\def\Curvesstack{\mathcal{C}\!\mathit{urves}}
\def\Polarizedstack{\mathcal{P}\!\mathit{olarized}}
\def\Complexesstack{\mathcal{C}\!\mathit{omplexes}}
% \Pic is the operator that assigns to X its picard group, usage \Pic(X)
% \Picardstack_{X/B} denotes the Picard stack of X over B
% \Picardfunctor_{X/B} denotes the Picard functor of X over B
\def\Pic{\mathop{\mathrm{Pic}}\nolimits}
\def\Picardstack{\mathcal{P}\!\mathit{ic}}
\def\Picardfunctor{\mathrm{Pic}}
\def\Deformationcategory{\mathcal{D}\!\mathit{ef}}

% --- appended 2026-10-12 (Chief): the project own definitions of the reference/slogan/history
% display environments, copied mechanically from the shared Stacks preamble lines 8-12
% (\usepackage{verbatim} plus the three \newenvironment lines, which the chapter preamble keeps
% verbatim from the project preamble). Nothing added, removed or redefined: these environments
% carry no mathematical content and the project itself typesets them as comments. Without these
% three lines a counted statement whose own text uses them fails with Environment slogan undefined.
\usepackage{verbatim}
\newenvironment{reference}{\comment}{\endcomment}
\newenvironment{slogan}{\comment}{\endcomment}
\newenvironment{history}{\comment}{\endcomment}
\usepackage{cleveref}
\crefname{originallemma}{original lemma}{original lemmas}
\Crefname{originallemma}{Original lemma}{Original lemmas}
\makeatletter
\newlabel{more-algebra-lemma-flat-descent-flat-at-primes-over}{{05LL}{source}{}{sourcepointer-1}{}}
\newlabel{more-algebra-lemma-flat-descent-flat-at-primes-over@cref}{{[originallemma][0][]05LL}{[1][0][]source}{}{}{}}
\makeatother
\crefname{originallemma}{original lemma}{original lemmas}
\Crefname{originallemma}{Original lemma}{Original lemmas}
\makeatletter
\newlabel{equation-flat-at-primes-over}{{052W}{source}{}{sourcepointer-2}{}}
\newlabel{equation-flat-at-primes-over@cref}{{[originallemma][0][]052W}{[1][0][]source}{}{}{}}
\makeatother
\crefname{originallemma}{original lemma}{original lemmas}
\Crefname{originallemma}{Original lemma}{Original lemmas}
\makeatletter
\newlabel{algebra-lemma-base-change-flat-up-down}{{00MQ}{source}{}{sourcepointer-3}{}}
\newlabel{algebra-lemma-base-change-flat-up-down@cref}{{[originallemma][0][]00MQ}{[1][0][]source}{}{}{}}
\makeatother
\title{Descent of the flat-at-primes-over-an-ideal condition along a surjective map on the closed sets}
\author{The Stacks Project Authors}
\date{stacks-project:more-algebra.tex:local-snapshot-20261012; GNU FDL 1.2 or later}
\begin{document}
\maketitle
\setcounter{subsection}{1}
\noindent\textit{Editorial note.} All delivered mathematical text is the author's own, copied from the pinned Stacks Project chapter file more-algebra.tex; no statement or proof step was rewritten, reordered or removed.\par\medskip
\noindent\textit{Editorial note.} The chapter file pinned-original-source.tex carries only \textbackslash{}input\{preamble\} on its first line, so no preamble text lives inside the chapter. The retained file sources/stacks-more-algebra/retained-preamble.tex is a mechanical, line-for-line copy of the project's single shared preamble as staged in the snapshot (Xy-pic setup plus theorem declarations and mathematical macros). Nothing was added, removed or redefined, and the excerpt serves only as preamble, never as a span coordinate.\par\medskip
\noindent\textit{Editorial note.} Both counted spans start at the \textbackslash{}begin line of their environment, so the statement and the author's proof are delivered whole; surrounding prose of the section is deliberately outside the counted spans and is carried as standard prerequisites instead.\par\medskip
\noindent\textit{Editorial note.} The citation fields (source\_id, source\_version, authors, source\_url, license.name, source\_title) are identical to those of every other card this lane has produced from this chapter, so the batch selector treats them as one source.\par\medskip
\noindent\textit{Allowed standard prerequisites (editorial).}\begin{itemize}
\item Flatness at a prime descends along a flat local base change. Algebra, Lemma \textbackslash{}ref\{algebra-lemma-base-change-flat-up-down\}, cited at proof lines 4135-4137. Carries flatness of the localised base changed module back to the localisation of the module over the base.
\item Surjectivity of the induced map on the closed sets defined by the ideals. assumption (1) of the lemma, used at proof lines 4125-4127. Produces the prime of the base change lying over the given prime.
\end{itemize}\medskip
\ifdefined\hypertarget\hypertarget{sourcepointer-1}{}\fi
\noindent\textit{Source reference.} Original lemma 05LL: focused-seeds/AGP/sources/stacks-snapshot-20261011/tags.txt, source lines 7240--7240. The lemma's own canonical tag, quoted so that readers can resolve the label inside the delivered text; it points at the counted statement itself and carries no proof content.\par\medskip
\ifdefined\hypertarget\hypertarget{sourcepointer-2}{}\fi
\noindent\textit{Source reference.} Original lemma 052W: focused-seeds/AGP/sources/stacks-snapshot-20261011/tags.txt, source lines 6584--6584. Navigation pointer to a statement the counted proof refers to but does not contain; nothing load-bearing is delegated to it.\par\medskip
\ifdefined\hypertarget\hypertarget{sourcepointer-3}{}\fi
\noindent\textit{Source reference.} Original lemma 00MQ: focused-seeds/AGP/sources/stacks-snapshot-20261011/tags.txt, source lines 832--832. Navigation pointer to a statement the counted proof refers to but does not contain; nothing load-bearing is delegated to it.\par\medskip
\begin{lemma}
\label{lemma-flat-descent-flat-at-primes-over}
Let $R \to S$ be a ring map.
Let $I \subset R$ be an ideal.
Let $M$ be an $S$-module.
Let $R \to R'$ be a ring map and $IR' \subset I' \subset R'$ an ideal
such that
\begin{enumerate}
\item the map $V(I') \to V(I)$ induced by
$\Spec(R') \to \Spec(R)$ is surjective, and
\item $R'_{\mathfrak p'}$ is flat over $R$ for all primes
$\mathfrak p' \in V(I')$.
\end{enumerate}
If (\ref{equation-flat-at-primes-over}) holds for
$(R' \to S \otimes_R R', I', M \otimes_R R')$, then
(\ref{equation-flat-at-primes-over}) holds for $(R \to S, I, M)$.
\end{lemma}

\begin{proof}
Assume (\ref{equation-flat-at-primes-over}) holds for
$(R' \to S \otimes_R R', IR', M \otimes_R R')$. Pick a prime
$IS \subset \mathfrak q \subset S$. Let $I \subset \mathfrak p \subset R$
be the corresponding prime of $R$. By assumption there exists
a prime $\mathfrak p' \in V(I')$ of $R'$ lying over $\mathfrak p$ and
$R_{\mathfrak p} \to R'_{\mathfrak p'}$ is flat. Choose a prime
$\overline{\mathfrak q}' \subset
\kappa(\mathfrak q) \otimes_{\kappa(\mathfrak p)} \kappa(\mathfrak p')$
which corresponds to a prime $\mathfrak q' \subset S \otimes_R R'$ which
lies over $\mathfrak q$ and over $\mathfrak p'$. Note that
$(S \otimes_R R')_{\mathfrak q'}$ is a localization of
$S_{\mathfrak q} \otimes_{R_{\mathfrak p}} R'_{\mathfrak p'}$.
By assumption the module $(M \otimes_R R')_{\mathfrak q'}$ is
flat over $R'_{\mathfrak p'}$. Hence
Algebra, Lemma \ref{algebra-lemma-base-change-flat-up-down}
implies that $M_{\mathfrak q}$ is flat over $R_{\mathfrak p}$
which is what we wanted to prove.
\end{proof}

\end{document}
